%%% Hlavní soubor. Zde se definují základní parametry a odkazuje se na ostatní části. %%%

%% Verze pro jednostranný tisk:
% Okraje: levý 40mm, pravý 25mm, horní a dolní 25mm
% (ale pozor, LaTeX si sám přidává 1in)
\documentclass[12pt,a4paper]{report}
\setlength\textwidth{145mm}
\setlength\textheight{247mm}
\setlength\oddsidemargin{15mm}
\setlength\evensidemargin{15mm}
\setlength\topmargin{0mm}
\setlength\headsep{0mm}
\setlength\headheight{0mm}
% \openright zařídí, aby následující text začínal na pravé straně knihy
\let\openright=\clearpage

%% Pokud tiskneme oboustranně:
% \documentclass[12pt,a4paper,twoside,openright]{report}
% \setlength\textwidth{145mm}
% \setlength\textheight{247mm}
% \setlength\oddsidemargin{15mm}
% \setlength\evensidemargin{0mm}
% \setlength\topmargin{0mm}
% \setlength\headsep{0mm}
% \setlength\headheight{0mm}
% \let\openright=\cleardoublepage

\usepackage[czech]{babel}
\usepackage[utf8]{inputenc}
\usepackage[IL2]{fontenc}
\usepackage[labelfont={small,bf},textfont={small,it}]{caption}


%% Ostatní balíčky
\usepackage{graphicx}
\usepackage{amsthm}

%% Balíček hyperref, kterým jdou vyrábět klikací odkazy v PDF,
%% ale hlavně ho používáme k uložení metadat do PDF (včetně obsahu).
%% POZOR, nezapomeňte vyplnit jméno práce a autora.
\usepackage[unicode]{hyperref}   % Musí být za všemi ostatními balíčky
\hypersetup{
    pdftitle={Vývoj hlasově ovládaných webových her pomocí CloudASR},
	pdfauthor={Jan Milota}
}


\usepackage[intoc]{nomencl}
\makenomenclature

%%% Drobné úpravy stylu

% Tato makra přesvědčují mírně ošklivým trikem LaTeX, aby hlavičky kapitol
% sázel příčetněji a nevynechával nad nimi spoustu místa. Směle ignorujte.
\makeatletter
\def\@makechapterhead#1{
  {\parindent \z@ \raggedright \normalfont
   \Huge\bfseries \thechapter. #1
   \par\nobreak
   \vskip 20\p@
}}
\def\@makeschapterhead#1{
  {\parindent \z@ \raggedright \normalfont
   \Huge\bfseries #1
   \par\nobreak
   \vskip 20\p@
}}
\makeatother

% Toto makro definuje kapitolu, která není očíslovaná, ale je uvedena v obsahu.
\def\chapwithtoc#1{
\chapter*{#1}
\addcontentsline{toc}{chapter}{#1}
}

\def\printnom#1{
\renewcommand{\nomname}{#1}
\printnomenclature[2cm]
}

\def\nom#1#2{#1\nomenclature{#1}{#2}}

\begin{document}

% Trochu volnější nastavení dělení slov, než je default.
\lefthyphenmin=2
\righthyphenmin=2

%%% Titulní strana práce

\pagestyle{empty}
\begin{center}

\large

Univerzita Karlova v Praze

\medskip

Matematicko-fyzikální fakulta

\vfill

{\bf\Large BAKALÁŘSKÁ PRÁCE}

\vfill

\centerline{\mbox{\includegraphics[width=60mm]{img/logo.png}}}

\vfill
\vspace{5mm}

{\LARGE Jan Milota}

\vspace{15mm}

% Název práce přesně podle zadání
{\LARGE\bfseries Vývoj hlasově ovládaných webových her pomocí CloudASR}

\vfill

% Název katedry nebo ústavu, kde byla práce oficiálně zadána
% (dle Organizační struktury MFF UK)
Ústav formální a aplikované lingvistiky

\vfill

\begin{tabular}{rl}

Vedoucí bakalářské práce: & Mgr. Ing. Filip Jurčíček, Ph.D. \\
\noalign{\vspace{2mm}}
Studijní program: & Informatika \\
\noalign{\vspace{2mm}}
Studijní obor: & Obecná  Informatika \\

\end{tabular}

\vfill

% Zde doplňte rok
Praha 2015

\end{center}

\newpage

%%% Následuje vevázaný list -- kopie podepsaného "Zadání bakalářské práce".
%%% Toto zadání NENÍ součástí elektronické verze práce, nescanovat.

%%% Na tomto místě mohou být napsána případná poděkování (vedoucímu práce,
%%% konzultantovi, tomu, kdo zapůjčil software, literaturu apod.)

\openright

\noindent
Děkuji panu doktorovi Jurčíčkovi za profesionální vedení této práce a za jeho neutuchající víru v mou schopnost samostatného vývoje.

\newpage

%%% Strana s čestným prohlášením k bakalářské práci

\vglue 0pt plus 1fill

\noindent
Prohlašuji, že jsem tuto bakalářskou práci vypracoval samostatně a výhradně
s~použitím citovaných pramenů, literatury a dalších odborných zdrojů.

\medskip\noindent
Beru na~vědomí, že se na moji práci vztahují práva a povinnosti vyplývající
ze zákona č. 121/2000 Sb., autorského zákona v~platném znění, zejména skutečnost,
že Univerzita Karlova v Praze má právo na~uzavření licenční smlouvy o~užití této
práce jako školního díla podle §60 odst. 1 autorského zákona.

\vspace{10mm}

\hbox{\hbox to 0.5\hsize{%
V ............. dne .............
\hss}\hbox to 0.5\hsize{%
Podpis autora
\hss}}

\vspace{20mm}
\newpage

%%% Povinná informační strana bakalářské práce

\vbox to 0.5\vsize{
\setlength\parindent{0mm}
\setlength\parskip{5mm}

Název práce:
Vývoj hlasově ovládaných webových her pomocí CloudASR
% přesně dle zadání

Autor:
Jan Milota

Katedra:  % Případně Ústav:
Ústav formální a aplikované lingvistiky
% dle Organizační struktury MFF UK

Vedoucí bakalářské práce:
Mgr. Ing. Filip Jurčíček Ph.D., Ústav formální a aplikované lingvistiky
% dle Organizační struktury MFF UK, případně plný název pracoviště mimo MFF UK

Abstrakt:
% abstrakt v rozsahu 80-200 slov; nejedná se však o opis zadání bakalářské práce
Cílem práce je navrhnout a vyvinout software pro výuku jazyků hrou za~použití webových technologií a~čerstvě vznikající CloudASR knihovny. Běžný uživatel provozuje 
interakci se svým prohlížečem skoro výhradně prostřednictvím myši a klávesnice. Díky softwaru, který tato práce reprezentuje, má nyní uživatel příležitost zabřednout do 
někdy ne~úplně populární výuky jazyka i za~pomoci svého hlasu. Což nabízí zmíněné výuce netušené možnosti, obzvláště stran uživatelské interaktivity. Důraz byl kladen 
na~uživatelskou přívětivost, grafickou fidelitu a~na~kompetitivní aspekt výuky, využívajíc Facebookovou integraci a~bodové hodnotící žebříčky.

Klíčová slova:
automatické rozpoznávání řeči, ASR, hry, web, HTML5, Javascript

\vss}\nobreak\vbox to 0.49\vsize{
\setlength\parindent{0mm}
\setlength\parskip{5mm}

Title:
% přesný překlad názvu práce v angličtině
Development of speech enabled web games using CloudASR

Author:
Jan Milota

Department:
Institute of Formal and Applied Linguistics
% dle Organizační struktury MFF UK v angličtině

Supervisor:
Mgr. Ing. Filip Jurčíček Ph.D., Institute of Formal and Applied Linguistics
% dle Organizační struktury MFF UK, případně plný název pracoviště
% mimo MFF UK v angličtině

Abstract:
% abstrakt v rozsahu 80-200 slov v angličtině; nejedná se však o překlad
% zadání bakalářské práce
The main goal of this thesis is to design and implement a piece of software for playful language learning, using web technologies and the fresh CloudASR library. A 
common user interacts with their web browser almost exclusively using a mouse and keyboard. Thanks to the software this thesis represents the user has an opportunity to delve into sometimes unpopular language learning process using his natural voice. This fact presents new and exciting possibilities, mainly regarding user 
interactivity. A lot of stress has been put to user friendliness, graphical fidelity and to the competitive aspect of language education, exploiting Facebook integration and point-scoring leader boards.

Keywords:
% 3 až 5 klíčových slov v angličtině
automatic speech recognition, ASR, games, web, HTML5, Javascript

\vss}

\newpage

%%% Strana s automaticky generovaným obsahem bakalářské práce. U matematických
%%% prací je přípustné, aby seznam tabulek a zkratek, existují-li, byl umístěn
%%% na začátku práce, místo na jejím konci.

\openright
\pagestyle{plain}
\setcounter{page}{1}
\tableofcontents

%%% Jednotlivé kapitoly práce jsou pro přehlednost uloženy v samostatných souborech
\chapwithtoc{Úvod}

Základním komunikačním kanálem mezilidské interakce je mluvená řeč. Ačkoli trend v poslední době spěje spíše k separaci této komunikace virtuálními médii, je mluvené slovo i nadále nenahraditelnou součástí socializace a výměny informací. Proč tedy nepřenést schopnost porozumět lidské mluvě i na stroje, jakožto právě na ona separující virtuální média? Současná technologie je k tomu jistě dostačující, výpočetní výkon lidstva roste po zběsilé křivce a technologické know-how po ještě zběsilejší.

Aby tento cíl mohl být dosažen, je kritické mít dobrý \nom{ASR}{Automatic Speech Recognition - automatické rozpoznávání řeči} (Automatic Speech Recognition) systém. Těchto se pohybuje v éteru mnoho, volně dosažitelných, rozličné kvality a spolehlivosti. Pro specifické využití v této práci byl zvolen systém 
\verb|CloudASR|\footnote{\url{https://github.com/UFAL-DSG/cloud-asr/}}, jehož autorem je kolega Ondřej Klejch.

\texttt{\nom{TrogASR}{Translating Online Game using Automatic Speech Recognition}} (Translating Online Game using Automatic Speech Recognition) je, jak možná název napovídá, online hra, zaměřující se na překlad slov a frází z~jazyka do jazyka za využití automatizovaného rozpoznávání řeči.

Cílem této práce je návrh a implementace graficky přitažlivé hry pro webové prohlížeče, která poskytne uživatelům další, neotřelou, možnost, jak si vyzkoušet a~zdokonalit své jazykové schopnosti. Tato hra pak bude sloužit jako kompetitivitou hnaný motor pro sběr verifikovaných dat pro použitou \verb|CloudASR| knihovnu.

Velký důraz je kladen na grafickou stránku věci; uživatelsky neatraktivní hra nemá valnou naději na šíření mezi uživateli samotnými samovolně. Navíc se na~ošklivou aplikaci nikdo nebude dívat po delší časový úsek, než je jedno sezení, a~už vůbec nikdo se k takovéto aplikaci nebude vracet. Proto bylo veškeré stylování a~malování a~navrhování layoutů a přechodů prováděno s výraznou pílí.

Dále je důraz kladen i na zmíněnou kompetitivitu. Každý má touhu vidět své jméno někde vysoko na žebříčku, zanechávaje své přátele na tomto žebříčku daleko pod sebou. Proto byla zapojena i integrace s Facebookovým API, jež~přináší kýženého výsledku poměrně nenásilnou formou. Na \uv{Like} a \uv{Login} tlačítko jsou uživatelé zvyklí a~nebojí se jej extenzivně používat.

Následující text popisuje a dokumentuje vývoj této aplikace a~myšlenkové pochody při něm stojící. První kapitola analyzuje e-learning jako způsob předávání informací. Druhá kapitola zmiňuje ASR teorii, popisuje některé postupy a základní paradigmata. Třetí kapitola seznámí čtenáře s využitými technologiemi a frameworky pro vývoj aplikace \verb|trogASR|. Kapitola čtvrtá nabídne čtenáři vhled do řešení a~implementace aplikace \verb|trogASR|, do použitých návrhových vzorů a do problematiky s nimi spojené. A konečně, kapitola pátá shrne celou práci.
\chapter{e-Learning jako způsob předávání informací}

\section{Kapitola}

\section{Kapitola}


\chapter{Automatické rozpoznávání řeči}

\section{Kapitola}

\section{Kapitola}


\chapter{Použité technologie}

V této kapitole se čtenář seznámí s technologiemi, frameworky a prostředími, využitými pro vývoj aplikace \verb|trogASR|. U každého z nich se navíc ve stručnosti dozví o díle zapojení daného nástroje do aplikace.

Všechny nástroje jsou buď {\sl open-source}, {\sl freeware}, nebo je pro ně volně dostupná akademická licence.

\section{Python}

\verb|Python|\footnote{\url{https://www.python.org/}} je {\sl open source} vysokoúrovňový skriptovací programovací jazyk. Nabízí širokou škálu programátorských paradigmat, jako například \nom{OOP}{Object Oriented Programming} (objektově orientované programování), imperativní, procedurální, nebo funkcionální přístup k programování.

Tento jazyk v osmdesátých letech dvacátého století navrhl Guido van Rossum. V roce 1989 započal vývoj. Od té doby popularita \verb|Python|u vzrostla natolik, že se stal základní součástí většiny distribucí systému Linux. Lze jej však velmi snadno instalovat i do jiných operačních systémů.

Hlavními návrhovými milníky a posléze i výhodami \verb|Python|u jsou jeho orientace na čistotu a efektivitu psaní kódu, na velkou modularitu a na výraznou uživatelskou přívětivost. Některými dokonce \verb|Python| bývá považován za nejvhodnější programovací jazyk pro začátečníky.

V aplikaci \verb|trogASR| bylo \verb|Python|u využito pro tvorbu {\sl back-end}u, tj. pro tvorbu prostředí aplikačního serveru za pomoci \nom{MVC}{Model-View-Controller} (Model-View-Controller) návrhového vzoru (více v sekci \ref{ssec:mvc}) a pro operování s databází.

\section{Flask}

\verb|Flask|\footnote{\url{http://flask.pocoo.org/}} je {\sl microframework} napsaný v \verb|Python|u, založený na knihovně \verb|Werkzeug|\footnote{\url{http://werkzeug.pocoo.org/}} a na templatovacím systému \verb|Jinja2|\footnote{\url{http://jinja.pocoo.org/}}.

Označení {\sl microframework} nese \verb|Flask| proto, že se jedná vskutku o framework minimalistický. Základní distribuce obsahuje pouze nástroje nezbytně nutné pro běh webové aplikace. Velký důraz je však kladen na modularitu tohoto frameworku; uživatel má tedy možnost doinstalovat si do svého prostředí právě ta rozšíření, která potřebuje.

Obrázek \ref{fig:flask_install} ukazuje, že instalace \verb|Flask|u je velice snadná. Stačí mít funkční prostředí \verb|Python|u a o zbytek se již může postarat nástroj \verb|easy_install| a \verb|pip|. Pro zprovoznění základní webové aplikace poté stačí už jen sedm řádek kódu.
\\
\begin{figure}[h]
	\verbatiminput{snip/flask.txt}
	\caption{Příklad instalace frameworku Flask na MS Windows, kódu modulu "hello"\;a spuštění webového serveru.}
	\label{fig:flask_install}
\end{figure}

Frameworku \verb|Flask| bylo v aplikaci \verb|trogASR| využito především proto, že tento framework je velmi jednoduchý na použití; naprosto oproštěn od přebytečných součástí, které by bylo třeba konfigurovat, ačkoli by je aplikace nevyužívala. Zároveň \verb|Flask| oplývá širokou komunitou uživatelů, schopných poradit s problémem, a exhaustivní dokumentací.

\subsection{Flask-Assets}

\begin{figure}[h]
	\verbatiminput{snip/flask_assets.txt}
	\caption{Příklad instalace balíčku Flask-Assets na MS Windows, zapojení do modulu "hello"\;a ukázka použití v Jinja2 HTML templatu.}
	\label{fig:flask_a_install}
\end{figure}

\verb|Flask-Assets|\footnote{\url{http://flask-assets.readthedocs.org/}} je rozšiřovací balíček pro \verb|Flask| framework, který řeší správu a minifikaci statických skriptů a stylů.

Jak ukazuje obrázek \ref{fig:flask_a_install}, instalace je opět triviální, používání balíčku pak velmi intuitivní a praktické.

Rozšíření \verb|Flask-Assets| lze také jednoduše konfigurovat. Například spouštíme-li aplikaci v {\sl debug} módu, můžeme knihovně říci, aby poskytovala nesloučené a neminifikované zdroje.

Tato knihovna je v aplikaci \verb|trogASR| využívána hlavně z důvodu pohodlnosti správy statických zdrojů.

\subsection{SQLite3}

Pro správu databáze byl zvolen {\sl light-weight} framework \verb|SQLite3|. Tento přichází již zahrnut jako součást instalace frameworku \verb|Flask|.

Aplikace \verb|trogASR| využívá tento framework kvůli jeho jednoduchosti. V budoucnosti je jednoduše možné jej nahradit některým z pokročilejších frameworků.

\section{JavaScript}

\verb|JavaScript| (\nom{JS}{JavaScript}) je multiplatformní, objektově orientovaný, interpretovaný skriptovací jazyk. Jeho autorem je Brendan Eich ze společnosti Netscape.

JS sice svou syntaxí patří do rodiny jazyků \verb|C/C++/Java|, s posledním zmíněným jazykem však krom slova \verb|Java| v názvu a podobné syntaxe nemá nic společného. Volba názvu má čistě marketingové důvody -- těžila právě z oblíbenosti jazyka \verb|Java|. Standardizovaným jménem JS je podle \nom{ISO}{International Organization for Standardization} (International Organization for Standardization) \verb|ECMAScript|.

JS je dnes již tak rozšířen, že jeho interpretaci podporuje drtivá většina současných webových prohlížečů. Nejčastějším využitím JS je totiž manipulace s WWW stránkou a ovládání všemožných prvků \nom{GUI}{Graphic User Interface} (Graphic User Interface).

Protože aplikace \verb|trogASR| je velmi dynamická a odesílat data při každé uživatelské akci na server by způsobovalo uživateli nemalou frustraci odezvou, je JS zhusta využíván. Pomocí JS je řešena naprostá většina animací, přechodů a dynamických prvků aplikace.

\subsection{jQuery}

\verb|jQuery|\footnote{\url{https://jquery.com/}} je malý, rychlý a možnostmi bohatý JS framework. Umožňuje uživateli zefektivnit proces psaní kódu, zjednodušit manipulaci s \nom{DOM}{Document Object Model} (Document Object Model) a obalit a zpříjemnit využití technologie \nom{AJAX}{Asynchronous JavaScript and XML} (Asynchronous JavaScript and XML).

Pro jeho jednoduchost, flexibilitu a ohromující uživatelskou komunitu byl \verb|jQuery| framework zvolen jako základní kámen celé aplikace \verb|trogASR|.

\subsection{transform.js}

\verb|transform.js|\footnote{\url{http://louisremi.github.io/jquery.transform.js/}} je minimalistický \verb|jQuery| plugin, umožňující animaci \nom{CSS}{Cascading Style Sheets} (Cascading Style Sheets -- kaskádové styly) vlastností {\sl transform}.

Jednoduchým zapojením a použitím získal tento plugin své místo v aplikaci \verb|trogASR|.

\section{HTML5}

\nom{HTML}{HyperText Markup Language} (HyperText Markup Language -- hypertextový značkovací jazyk) verze 5 je nejnovější iterací HTML jazyka. Přináší oproti svému předchůdci, verzi 4.01, nové sémantické prvky, vyřazuje zastaralé konstrukty z jazyka a opravuje řadu chyb.

HTML5 jazyka je v aplikaci \verb|trogASR| využito pouze poskrovnu; je nasazen pouze jako rámec aplikace, do kterého se většina obsahu generuje dynamicky s použitím JS.

Zároveň některých nových možností HTML5 využívá knihovna \verb|CloudASR| (jmenovitě API pro záznam zvuku).

\section{CSS3}

CSS technologie umožňuje programátorům a návrhářům oddělit strukturu a obsah HTML stránky od jejího vzhledu. CSS jazyk popisuje vztahy mezi elementy ve stránce a způsoby zobrazení. Díky tomuto popisu je možno vytvořit dobře znovupoužitelné nastavení, které je transparentně přenositelné mezi prvky nejen v jedné stránce, ale i mezi prohlížeči a platformami. Autorem konceptu CSS byl Håkon Wium Lie a standardizační organizace \verb|W3C|\footnote{\url{www.w3.org/}}.

Kaskádové styly verze 3 přináší na trh přehršel nových, užitečných konceptů. Mezi tyto patří například příchod {\sl @media} pravidla pro zachycení a podmínění stylů podle využitého zobrazovacího zařízení, nebo zavedení {\sl transform} a {\sl transition} vlastností elementů pro usnadnění pozicování a animace.

Aplikace \verb|trogASR| využívá CSS3 velmi extenzivně. CSS pravidla definují vzhled všech prvků, které uživatel vidí, a dokonce i některých animovaných přechodů.

\section{CloudASR}

\verb|CloudASR|\footnote{\url{http://www.cloudasr.com/}} je softwarová platforma a veřejná ASR služba. Za jejím vznikem stojí kolega Ondřej Klejch, který tuto platformu vyvinul v rámci své diplomové práce na MFF.

Hlavními rysy \verb|CloudASR| platformy jsou především 

\begin{itemize}
\item automatické škálování podle zatížení
\item jednoduchost nasazení a provozu
\item cloudová nátura
\item {\sl state-of-the-art} inkrementální rozpoznávání řeči
\item kompatibilita s {\sl Google Speech API}
\item rozšiřitelnost o nové ASR enginy
\end{itemize}

\verb|CloudASR| také poskytuje JS knihovnu, jejíž zapojení si aplikace \verb|trogASR| z titulu práce žádala. Této knihovny je využito právě pro transformaci hlasového vstupu na výstup textový. Její zavedení demonstruje obrázek \ref{fig:cloudasr}.

Obrázek \ref{fig:cloudasr_final} pak ukazuje, jak vypadají \nom{JSON}{JavaScript Object Notation} (JavaScript Object Notation) objekty, které JS knihovna dodává na konci samotného rozpoznávání.
\\
\begin{figure}[h]
	\verbatiminput{snip/cloudasr.txt}
	\caption{Ukázka zapojení CloudASR JS knihovny.}
	\label{fig:cloudasr}
\end{figure}

\begin{figure}[h]
	\verbatiminput{snip/cloudasr_final.txt}
	\caption{Příklad finálního výsledku rozpoznávání.}
	\label{fig:cloudasr_final}
\end{figure}

\section{PyCharm}

\verb|PyCharm|\footnote{\url{https://www.jetbrains.com/pycharm/}} je vývojové prostředí, vyvíjené a udržované společností \verb|JetBrains|. Je určeno převážně pro vývoj aplikací v jazyce \verb|Python|, oplývá však i velmi sofistikovaným prostředím pro vývoj v JS.

JS podpora sahá až na úroveň nativního zapojení všemožných frameworků, jako jsou například právě \verb|jQuery|, ale i \verb|ExtJS|, nebo \verb|Node.js|.

Obecně je práce v tomto vývojovém nástroji velmi příjemná. Kontextová inteligentní nápověda, možnosti refactoringu a projektově-přehledové nástroje hledají na dnešním trhu kompetici jen těžce.

\section{Několik slov o webových prohlížečích}

V současné době je trh s webovými prohlížeči poměrně rozmanitý. Běžný uživatel má na výběr nepřeberné množství různých řešení, více i méně rozšířených a použitelných.

Bohužel, ani po několika letech existence norem HTML5 a CSS3 nejsou tyto v naprosté většině zmíněných prohlížečů podporovány, nebo jsou podporovány jen částečně, nebo chybně.

Proto aplikace \verb|trogASR| necílí na multiplatformnost. Bude fungovat pouze v prohlížečích, které dostatečně podporují čisté HTML5 a CSS3. Těmito jsou k dnešnímu dni pouze prohlížeče \verb|Google Chrome| (verze 41) a \verb|Opera Next| (verze 29).
\chapter{Implementace}
\label{chap:solution}

V této kapitole se čtenář dočte o implementaci aplikace \verb|trogASR|. Bude obeznámen s pilíři aplikace, nahlédne na vývoj aplikace jako na řešení komplexní, i po částech.

Nejprve se čtenář dočte o struktuře aplikace --- o modularitě, designu, znovupoužití kódu --- při čemž se seznámí s myšlenkovými pochody, stojícími za návrhem implementace.

Dále se čtenář dozví o interoperabilitě s Facebook \nom{API}{Application Programming Interface} (Application Programming Interface -- aplikační programátorské rozhraní) a o důležité součásti --- internacionalizaci aplikace.

V závěru je pak věnován prostor úhrnu některých použitých návrhových vzorů, principu jejich fungování a projekci do aplikace.

\section{Modulární kostra}

Aplikace \verb|trogASR|, jakkoli neobsahuje žádné extrémně složité algoritmické řešení, byla navrhována a vyvíjena s ohledem na rozšiřitelnost a modularitu. 
\\
\begin{figure}[h]
	\verbatiminput{snip/app.txt}
	\caption{Ukázka části kódu hlavní komponenty.}
	\label{fig:app}
\end{figure}

\begin{figure}[h]
	\verbatiminput{snip/comp.txt}
	\caption{Ukázka části internacionalizační komponenty.}
	\label{fig:comp}
\end{figure}

Každá komponenta systému musí být registrována v komponentě hlavní, aplikační, a musí splňovat základní implementační požadavky, aby vůbec systémovou komponentou být mohla.

Jelikož v JS neexistuje princip {\sl interface}, muselo být se smluvním typováním nakládáno jinak. Zvoleným přístupem je návrhový vzor {\sl duck-typing} (viz sekce \ref{ssec:ducktyping}).

Prvním krokem je běh JS skriptu, kdy prohlížeč skript interpretuje, jak je vložen do stránky. V tuto chvíli se na prvním místě v pořadí nachází skript inicializující hlavní aplikační komponentu. V této se pak nachází mechanismus pro registraci ostatních komponent.

Každá registrovaná komponenta musí splňovat několik základních vlastností -- musí zpřístupňovat metodu \verb|initComponent()| a metodu \verb|reset()|. První zmíněná metoda je volána na každé komponentě v handleru \verb|window.onload|, zatímco druhá zmíněná metoda je volána při nutnosti znovuvytvoření statických prvků komponenty (například při změně lokále, či při změně velikosti okna). Část implementace aplikační komponenty demonstruje obrázek \ref{fig:app}. Část implementace běžné komponenty pak zobrazuje obrázek \ref{fig:comp}.

\section{Templatovací systém}
\label{sec:template}

Pro pohodlnost používání a zvětšení znovupoužitelnosti kódu byl do aplikace \verb|trogASR| implementován systém na vytváření a processing HTML templatů. Tohoto systému celá aplikace využívá ve chvíli, potřebuje-li inicializovat jakoukoli komponentu vícekrát (položky menu, tlačítka, ...).

Každý takovýto template má formu HTML kódu, mixovaného s možným JS kódem; vymezeným konstruktem \verb|<% <JS_code> %>|, nebo \verb|<%= <JS_variable> %>|.
\\
\begin{figure}[h]
	\verbatiminput{snip/template.txt}
	\caption{Ukázka "facebook widget"\;templatu a jeho použití.}
	\label{fig:template}
\end{figure}

Všechny tyto templaty jsou uzavřeny v \verb|<script type="text/html"></script>| elementech, kterým interpretující prohlížeč nepřikládá žádnou zvláštní váhu, pouze je vkládá do stránky. Příklad demonstruje obrázek \ref{fig:template}.

\section{Sprite factory a obrázkový font}

Pro účely zvýšení grafické fidelity byla navržena oku lahodící barevná paleta a ji využívající obrázkový font. Tento font je používán pro všechny texty, kde to má smysl, a slova v tomto fontu jsou generována dynamicky pomocí komponenty {\sl sprite factory}.

{\sl Sprite factory} komponenta se stará o vytváření dynamických grafických prvků podle současného nastavení jazyka a velikosti zobrazovacího média. Pro tento proces využívá taktéž templatovací systém (sekce \ref{sec:template}).
\\
\begin{figure}[h]
	\includegraphics[width=100mm]{img/sprite.png}
	\caption{Položka menu "Play"/"Hrát"\;pro anglické, respektive české jazykové nastavení.}
	\label{fig:sprite}
\end{figure}

Slova obrázkového fontu jsou tvořena ze {\sl spritů} jednotlivých písmen, uzavřených do obalovacího \verb|<div>| elementu. Velikost písmen a slov je dána responzivním nastavením aplikace (více v sekci \ref{sec:responsive}). Každé písmeno je reprezentováno vlastním obrázkem, pro nějž existuje CSS styl, který jej zapojí jako pozadí příslušného elementu. Vzhled takto vygenerovaného dynamického slova ukazuje obrázek \ref{fig:sprite}.

\section{Responzivní design}
\label{sec:responsive}

Protože aplikace \verb|trogASR| se má líbit a být uživatelsky oblíbená, vyvstala nutnost možnosti aplikaci zobrazovat na co možná nejširším spektru zobrazovacích médií. Ačkoli aplikace podporuje jenom specifickou podsekci webových prohlížečů, podporované prohlížeče lze nalézt i na jiných platformách, než jen PC (\nom{PC}{Personal Computer} -- Personal Computer) --- například mobilní zařízení s operačním systémem Android. Zároveň je třeba vzít v potaz i fakt, že zobrazovací média uživatelů běžných PC se liší, někdy i výrazně.

Proto je aplikace \verb|trogASR| navržena čistě dynamicky. Veškeré stylování je činěno pomocí procentuálních poměrů velikostí elementů vůči jejich nadřazeným elementům. 

V místech, kde je procentuální vyjádření nemožné (šířky okrajů, velikosti {\sl box-shadow} vlastností), jsou nasazeny jednotky \verb|vw| a \verb|vh|. Tyto jednotky vyjadřují jedno procento šířky zobrazované plochy (\verb|1vw|), respektive jedno procento její výšky (\verb|1vh|).
\\
\begin{figure}[h]
	\verbatiminput{snip/responsive.txt}
	\caption{Ukázka propagace dynamické velikosti písmen obrázkového fontu.}
	\label{fig:responsive}
\end{figure}

Dynamická velikost písmen obrázkového fontu je zařízena pomocí skrytého \verb|<div>| elementu, který je nastylován na tři procenta šířky zobrazovacího média. V inicializaci komponenty {\sl sprite factory} je posléze tato skutečnost zaznamenána a aplikována na jednotlivá písmena jako jejich výška a šířka (písmena jsou vždy čtvercová). Rozměr písmen a slov je dále možno upravovat posláním multiplikačního parametru velikosti do příslušných metod. Tento proces demonstruje obrázek \ref{fig:responsive}.

Obecně je doporučený minimální poměr stran 1:1, maximální pak 2:1. Optimální poměry jsou 16:9, nebo 4:3. Doporučené minimální rozlišení je {\sl 720p \nom{HD}{High Definition}} (High Definition), optimální {\sl 1080p \nom{FHD}{Full HD}} (Full HD), maximální {\sl 4k \nom{UHD}{Ultra HD}} (Ultra HD).

\section{Facebook API}

Hra \verb|trogASR| by také měla být kompetitivní. Proto byla zvolena integrace s facebookovým API. Tato zajišťuje, že chce-li uživatel soupeřit s ostatními, má jednoduchou možnost, jak se do aplikace přihlásit.

V případě, že uživatel přihlašovací možnosti nevyužije, má i tak možnost hru hrát, nemůže se však umístit na žebříčcích. Aplikace jako taková na uživatelovu \uv{zeď} žádné příspěvky neodesílá --- uživatel může tuto volbu učinit sám pomocí \uv{Like}, nebo \uv{Share} tlačítek, které se mu taktéž implicitně nabízí. Není proto důvod, aby se uživatel do hry nepřihlásil. Navíc, na facebooková tlačítka jsou lidé poměrně zvyklí a nebojí se jich použít.

Pro zapojení facebookových tlačítek do webové stránky nabízí Facebook poměrně jednoduché a přímočaré řešení\footnote{\url{https://developers.facebook.com/docs/javascript}}. Stačí do stránky umístit \verb|<div>| element, přiřadit mu příslušnou CSS třídu, zadefinovat příslušné {\sl data} atributy a natáhnout do stránky skript, který všechny tyto elementy přepíše elementy {\sl iframe}. Do těchto pak načte svá data a zobrazí je. Příkladem nastavení facebookových {\sl widgetů} budiž obrázek \ref{fig:template}.

Velkou nevýhodou facebookových {\sl widgetů} je ale poměrně malá konfigurovatelnost. Obzvlášť co se týče nastavení velikostí tlačítek, neposkytuje API --- krom předdefinovaných {\sl small}, {\sl medium}, {\sl large} rozměrů --- žádné rozšiřující možnosti. Což v důsledku působí problém s ohledem na responzivní design aplikace. Je-li zobrazovací médium moc malé, tlačítka začnou \uv{vytékat} z přiřazených kontejnerů. Je-li naopak moc velké, tlačítka jsou titěrně malá.

\section{Internacionalizace}

Anglická zkratka \nom{I18N}{Internationalization} (Internationalization) je běžně používána pro označení částí systému, které jsou schopny přizpůsobit své chování (nebo toto přizpůsobení zajistit) dle volby uživatelského lokále. Aplikace \verb|trogASR| tento koncept přejímá a integruje.

Každý zobrazovaný text, který má být internacionalizován, musí mít svůj lokalizační klíč. Pro každý klíč pak musí existovat překlady do všech podporovaných jazyků. Na překlady textů je možno přistupovat buďto skrze statickou proměnnou \verb|t| modulu I18N, kam se příslušné texty uloží vždy při změně lokále, nebo přes dynamickou \verb|getText([locale])| metodu, která na vstupu požaduje lokále, pro nějž se má text získat. V případě vynechání tohoto parametru metoda na klíči vrátí text pro současně nastavené lokále. Takže například text pro tlačítko {\sl play} lze získat buďto ze statického kontextu jako \verb|APP.I18N.t.MENU_PLAY|, nebo dynamicky jako \verb|APP.I18N.keys.MENU_PLAY.getText([locale])|.

V současné době jsou v základu podporovány dva zobrazovací jazyky --- angličtina a čeština. Nastavení jazyka přepne kompletně celou aplikaci a vynutí tedy její reset.

Zároveň je možno poměrně jednoduše v několika málo krocích přidat další zobrazovací jazyk. Je pro to potřeba:

\begin{itemize}
\item zavést nové lokále do souboru \verb|i18n.js| do příslušného výčtu s příslušnými zkratkami ({\sl en, cs, klingon, ...})
\item přidat překlad pro každý lokalizační klíč v témže souboru
\item poskytnout styly pro tlačítko na výběr nového lokále (viz obrázek \ref{fig:locale_style})
\item pokud nový jazyk má ve své abecedě nějaké speciální znaky, je nutno poskytnout i tyto (viz obrázek \ref{fig:locale_letter})
\end{itemize}

O zbytek se postará jádro modulu I18N.
\\
\begin{figure}[h]
	\verbatiminput{snip/locale_style.txt}
	\caption{Příklad přidání stylu pro ikonu výběru jazyka \uv{klingon}.}
	\label{fig:locale_style}
\end{figure}

\begin{figure}[h]
	\verbatiminput{snip/locale_letter.txt}
	\caption{Příklad přidání specifického znaku do abecedy nového jazyka a nastavení nového prvku obrázkového fontu.}
	\label{fig:locale_letter}
\end{figure}

\subsection{Zdrojové a cílové jazyky překladu}

Jelikož se hra \verb|trogASR| zaměřuje na překlad slov a frází z jazyka do jazyka, lze pochopitelně nastavit i dvojice zdrojový a cílový jazyk. Zdrojový jazyk je ten, ve kterém se uživateli budou slova a fráze zobrazovat, cílový naopak ten, ve kterém uživatel slova a fráze musí vyřknout.

Platforma \verb|CloudASR| v současnosti rozumí pouze češtině a angličtině, proto cílovým jazykem musí být vždy jeden z nich. Nic však nebrání v zavedení nového zdrojového jazyka. Pro jeho přidání je potřeba:

\begin{itemize}
\item pokud zdrojový jazyk obsahuje speciální znaky, nutno přidat znak do obrázkového fontu (viz obrázek \ref{fig:locale_letter})
\item zaregistrovat dvojici zdrojový - cílový jazyk v příslušném výčtu v modulu {\sl options} (viz obrázek \ref{fig:locale_options})
\item přidat lokalizační klíč a překlady pro zobrazení výběru dvojice jazyků v nastavení do modulu I18N
\item poskytnout slovník v \nom{CSV}{Comma Separated Values} (Comma Separated Values -- hodnoty oddělené čárkami) formátu, kde první slovo na každé řádce je slovo ve zdrojovém jazyce, ostatní hodnoty na řádce jsou významy slova v jazyce cílovém (\verb|vergh;car;vehicle;starship|)
\item naimportovat tento slovník do databáze za využití připravené funkcionality (viz obrázek \ref{fig:locale_import})
\end{itemize}

\begin{figure}[h]
	\verbatiminput{snip/locale_options.txt}
	\caption{Příklad zavedení dvojice \uv{klingon} - \uv{en} do modulu {\sl options}.}
	\label{fig:locale_options}
\end{figure}

\begin{figure}[h]
	\verbatiminput{snip/locale_import.txt}
	\caption{Příklad importu slovníku \uv{klingon} $\rightarrow$ \uv{en}.}
	\label{fig:locale_import}
\end{figure}

\section{Návrhové vzory}

V aplikaci \verb|trogASR| je využita celá škála návrhových vzorů. Převážně byly tyto zapojovány pro zpřehlednění hektického a chaotického světa JS. Některé samospádem vyplývají z konceptu zpracování, některé jeden ze druhého. Každý je však unikátní a důležitý svým vlastním způsobem. Následuje výčet a popis některých z nich.

\subsection{Duck-typing}
\label{ssec:ducktyping}

\epigraph{\sl\uv{Pokud vidím ptáka, který chodí jako kachna, plave jako kachna a kváká jako kachna, tak o tomto ptáku tvrdím, že je to kachna.}}{--- James Whitcomb Riley}

{\sl Duck-typing} (\nom{DT}{Duck-Typing}) je metoda dynamického typování objektů, při které zkoumáme objekt z hlediska jeho vlastností a ne z hlediska implementovaných rozhraní a předků. Tato metoda se chová podobně, jako kdybychom nějaké rozhraní vskutku implementovali, nemusíme tak však činit; stačí nám implementovat metody, které rozhraní udává.

Zmíněný epigraf tedy není stoprocentně spojitelný s principem DT. Po drobné úpravě raději tvrdíme, že \uv{pokud vidím něco, co umí chodit, kvákat a plavat, pak věřím tomu, že je to kachna}.

Důležité je pozorování, že stačí vidět \uv{něco}, a že poté \uv{věříme}, že to je daného typu. Kdybychom pomocí DT například hledali mezi dostupnými třídu \verb|Enemy| jako implementora metody \verb|engage()|, mohli bychom se velmi divit, kdybychom jako první našli třídu \verb|Girlfriend|.

Jelikož JS je jazyk, který nezná koncept rozhraní, DT je jediná možnost, jak typovost zajistit. \verb|TrogASR| aplikace tuto možnost využívá při skládání modulárního systému.

\subsection{Module pattern}

{\sl Module pattern} je vzor určený pro organizaci JS aplikace. Zapouzdřuje vznik jednotlivých modulů do volání anonymní funkce, kde modul může vytvořit svou vnitřní strukturu. Ven pak vysouvá pouze API, které je s to poskytnout.

Tímto postupem třímáme schopnost dostat do lokálního rámce globální proměnné, které by JS interpret jinak musel složitěji vyhledávat v celém kódu. Jednoduše je pošleme jako parametry tovární funkce. Máme takto postaráno o přehlednější strukturu a jsme vyzbrojeni alespoň částečnou ochranou proti injekci potenciálně škodlivého kódu.

Aplikace \verb|trogASR| --- její prezentační část --- je na tomto vzoru celá postavena.

\subsection{Scope cache pattern}

Jelikož JS je jazyk s funkcionálními prvky, kde proměnná běžně může být typu funkce, musí znát princip {\sl scope} (rámec). Vnořené funkce mají vždy přehled o svém rámci a o v něm dostupných proměnných. Proto můžeme elegantně vytvářet statické funkce, které dokážou mít přístup do dynamického obsahu je obklopujícího.

Tím, že zavádíme a vydáváme funkce, které s sebou nesou informaci o stavu svého tvůrce, lze obejít jindy nevyhnutelnou nutnost statického předávání konfigurační informace, nebo nutnost posílání extrémního množství funkčních argumentů.

Tohoto vzoru aplikace \verb|trogASR| využívá poměrně nemálo; všude, kde lze. Exemplárním případem budiž generování statických \verb|getText()| metod na lokalizačních klíčích v I18N modulu, které se vždy při svém volání zařizují podle dynamické konfigurace modulu.

\subsection{Singleton pattern}

Název {\sl singleton pattern} vychází z matematické definice, která jako {\sl singleton} označuje množinu s právě jedním prvkem. V informatice se termínem {\sl singleton} označuje třída, jež se instancuje vždy maximálně jednou.

Ve světě JS je tento vzor poměrně obvyklý. Často jsou zde využívány globální proměnné, kterými jsou objekty zrovna typu {\sl singleton}. Aplikace \verb|trogASR| v tomto ohledu není výjimkou a vzor využívá též. Aplikační moduly jsou všechny tohoto typu.

\subsection{Currying}

{\sl Currying}, česky také curryfikace, se jako termín poprvé objevuje v jazyce \verb|Haskell| a získal svůj název (stejně jako jazyk) podle amerického matematika a logika Haskella Curryho.

Currifikace demonstruje základní paradigma funkcionálního programování. Jedná se o částečnou aplikaci funkce na parametry, kde výsledkem je opět funkce, která přijímá parametry dosud neaplikované. Obrázek \ref{fig:curry} ukazuje kontrast curryfikace v Haskellu a v JS.

\begin{figure}[h]
	\verbatiminput{snip/curry.txt}
	\caption{Příklad curryfikace v Haskellu a v JS.}
	\label{fig:curry}
\end{figure}

Curryfikace je v aplikaci \verb|trogASR| využito při zpracování templatů, kdy se nejprve vyrábí funkce, která template podle vstupu částečně upraví, a teprve poté se tato aplikuje na příchozí data.

\subsection{MVC pattern}
\label{ssec:mvc}

MVC (Model-View-Controller) je jediný návrhový vzor, který byl využit v aplikaci \verb|trogASR| na {\sl back-end}u (serverové části). Není však pravidlem, že by MVC vzor měl být pouze tam. Například {\sl front-end} framework \verb|ExtJS| od verze 5.0 přináší možnost nasazení MVC i na prezentační vrstvě.

MVC vzor diktuje rozdělení aplikace na tři části --- datový {\sl model}, uživatelské rozhraní ({\sl view}) a řídící logiku ({\sl controller}) --- tak, aby modifikace kterékoli této části měla minimální vliv na části ostatní. Vpravdě se nejedná ani tak o vzor návrhový, jako spíše o vzor architektonický (někdy také agregační). Týká se totiž výrazněji architektury aplikace, než-li klasické vzory návrhové.

\verb|Flask| {\sl microframework} tento vzor implicitně podporuje. Proto bylo jeho nasazení více, než žádoucí. Vzor je navíc podpořen specifickou strukturalizací zdrojového kódu.

% Ukázka použití některých konstrukcí LateXu (odkomentujte, chcete-li)
%%% Ukázka použití některých konstrukcí LaTeXu

\subsection{Ukázka \LaTeX{}u}
\label{ssec:ukazka}

V~této krátké části ukážeme použití několika základních konstrukcí \LaTeX{}u,
které by se vám mohly při psaní práce hodit.

Třeba odrážky:

\begin{itemize}
\item Logo Matfyzu vidíme na obrázku~\ref{fig:mff}.
\item Tato subsekce má číslo~\ref{ssec:ukazka}.
\item Odkaz na literaturu
\end{itemize}

Druhy pomlček:
červeno-černý (krátká),
strana 16--22 (střední),
$45-44$ (minus),
a~toto je --- jak se asi dalo čekat --- vložená věta ohraničená dlouhými pomlčkami.
(Všimněte si, že jsme za \verb|a| napsali vlnovku místo mezery: to aby se
tam nemohl rozdělit řádek.)

% Makro na české uvozovky (novější verze LaTeXu ho už mají zabudované)
%\newcommand{\uv}[1]{\quotedblbase #1\textquotedblleft}
\uv{České uvozovky.}

\newtheorem{theorem}{Věta}
\newtheorem*{define}{Definice}	% Definice nečíslujeme, proto "*"

\begin{define}
{\sl Strom} je souvislý graf bez kružnic.
\end{define}

\begin{theorem}
Tato věta neplatí.
\end{theorem}

\begin{proof}
Neplatné věty nemají důkaz.
\end{proof}

\begin{figure}
	\centering
	\includegraphics[width=30mm]{img/logo.png}
	\caption{Logo MFF UK}
	\label{fig:mff}
\end{figure}


\chapter*{Závěr}
\addcontentsline{toc}{chapter}{Závěr}


%%% Seznam použité literatury
%%% Seznam použité literatury je zpracován podle platných standardů. Povinnou citační
%%% normou pro bakalářskou práci je ISO 690. Jména časopisů lze uvádět zkráceně, ale jen
%%% v kodifikované podobě. Všechny použité zdroje a prameny musí být řádně citovány.

\def\bibname{Seznam použité literatury}
\begin{thebibliography}{99}
\addcontentsline{toc}{chapter}{\bibname}

\bibitem{lamport94}
  {\sc Lamport,} Leslie.
  \emph{\LaTeX: A Document Preparation System}.
  2. vydání.
  Massachusetts: Addison Wesley, 1994.
  ISBN 0-201-52983-1.

\end{thebibliography}


%%% Tabulky v bakalářské práci, existují-li.
\chapwithtoc{Seznam tabulek}

%%% Použité zkratky v bakalářské práci, existují-li, včetně jejich vysvětlení.
\printnom{Seznam použitých zkratek}

%%% Přílohy k bakalářské práci, existují-li (různé dodatky jako výpisy programů,
%%% diagramy apod.). Každá příloha musí být alespoň jednou odkazována z vlastního
%%% textu práce. Přílohy se číslují.
\chapwithtoc{Přílohy}

\openright
\end{document}
